\section*{अध्याय \ref{ch:b2:climate-variability} --- जलवायु परिवर्तिता}

\begin{solution}{exo:b2:climate-variability:1}
(क) 2026 में $29.27 - 26.71 = \qty{2.56}{\celsius}$; 1997 में $28.79 - 26.73 = \qty{2.06}{\celsius}$।
(ख) सितंबर 2026 दोनों में आगे था: अधिक गर्म माह, \qty{0.48}{\celsius} से, और बड़ी विसंगति,
\qty{0.50}{\celsius} से। (ग) हर वर्ष की तुलना उस पर केंद्रित 30-वर्षीय आधार अवधि से की जाती है, जो
हर पाँच वर्ष में खिसकाई जाती है; 1980 और 1990 के दशकों का सितंबर \omterm{def:b2:climate-variability:modes}{जलवायु-माध्य}
2010 और 2020 के दशकों के \omterm{def:b2:climate-variability:modes}{जलवायु-माध्य} से थोड़ा भिन्न है।
\end{solution}

\begin{solution}{exo:b2:climate-variability:2}
(क) JJA: $(0.79 + 0.95 + 1.28)/3 = 1.007$; JAS: $(0.95 + 1.28 + 1.55)/3 = 1.260$; ASO: $(1.28 +
1.55 + 1.68)/3 = 1.503$; SON: $(1.55 + 1.68 + 2.01)/3 = 1.747$; OND: $(1.68 + 2.01 + 2.02)/3 =
1.903$, सभी \unit{\celsius} में। (ख) एक दशमलव तक पूर्णांकित: 1.0, 1.3, 1.5, 1.7, 1.9, यानी
तालिका के मान। (ग) ASO 2023 में, 1.5 पर।
\end{solution}

\begin{solution}{exo:b2:climate-variability:3}
(क) सूचकांक ताहिती के दाब में से डार्विन का दाब घटाकर उसके विचलन का अनुसरण करता है: ताहिती पर औसत से
नीचे और डार्विन पर ऊपर होने से वह ऋणात्मक है। (ख) एल~नीनो। (ग) दुर्बल पूर्वी व्यापारिक पवनें, पश्चिम में
सतह के निकट पश्चिमी पवनें भी, और इंडोनेशिया के ऊपर सामान्य से कम वर्षा, क्योंकि
गहरा संवहन पूर्व की ओर खिसक गया है।
\end{solution}

\begin{solution}{exo:b2:climate-variability:4}
(क) $804.1/880.6 = \qty{91.3}{\percent}$ औसत का: 90 से \qty{110}{\percent} के भीतर, एक
सामान्य ऋतु, यद्यपि उसके निचले भाग में। (ख) ऋतु का योग $556 + 659 + 427 + 252 =
\qty{1894}{mm}$ है; जुलाई उसका $659/1894 = \qty{34.8}{\percent}$ लाता है, एक-तिहाई से अधिक।
\end{solution}

\begin{solution}{exo:b2:climate-variability:5}
(क) $\tau = 1025 \times 3992 \times 50/16.1 = \qty{1.271e7}{s} = \qty{147.1}{\day}$, \qty{50}{m} के लिए,
और उसका दोगुना, \qty{294.2}{\day}, \qty{100}{m} के लिए। (ख) $\mathrm{e}^{-30.44/147.1} =
0.813$ और $\mathrm{e}^{-30.44/294.2} = 0.902$। (ग) अवमंदन काल $h$ के समानुपाती है: अधिक
उथली मिश्रित परत उसी प्रवाह के लिए कम ऊष्मा संचित करती है और विसंगति को जल्दी भूलती है, इसलिए
महासागर की स्मृति, और उससे मिलने वाली पूर्वानुमेयता, घटती है।
\end{solution}

\begin{solution}{exo:b2:climate-variability:6}
(क) तीन महीने 91.3 दिन हैं: $1.2 \times \mathrm{e}^{-91.3/147.1} = \qty{0.65}{K}$; छह महीने
182.6 दिन हैं: $1.2 \times \mathrm{e}^{-182.6/147.1} = \qty{0.35}{K}$। (ख) विसंगति
धीरे क्षीण होती है, इसलिए वह अगले महीने की विसंगति का शून्य से बेहतर अनुमान है। पर यादृच्छिक प्रणोदन
हर दिन नई विसंगतियाँ जोड़ता है, जो पूर्वानुमान के समय अज्ञात होती हैं; जैसे-जैसे आरंभिक विसंगति क्षीण होती है,
उनका अंश बढ़ता जाता है, और पूर्वानुमान \omterm{def:b2:climate-variability:modes}{जलवायु-माध्य} की ओर बहता है, जहाँ वह अब
सहायक नहीं रहता।
\end{solution}

\begin{solution}{exo:b2:climate-variability:7}
(क) $160 \times 86\,400 = \qty{1.38e7}{J}$, प्रति वर्ग मीटर और प्रति दिन \qty{13.8}{MJ}। (ख) परत
प्रति वर्ग मीटर $(1000 - 850) \times 100/9.81 = \qty{1529}{kg}$ वायु रखती है, जिसकी ऊष्मा
धारिता $1529 \times 1004.5 = \qty{1.536e6}{J.m^{-2}.K^{-1}}$ है; वह $1.38 \times
10^7/1.536 \times 10^6 = \qty{9.0}{K}$ प्रति दिन गर्म होगी। (ग) परत अवरक्त में ऊष्मा विकिरित करती है,
अंतरिक्ष की ओर और ऊपर की वायु की ओर; और गर्म वायु ऊपर उठकर अपने ऊपर की परत से मिलती है, या
पवनों द्वारा दूर ले जाई जाती है, जिनमें वह ठंडी वायु का अंतर्वाह भी है जिसे ऊष्मीय निम्न दाब स्वयं खींचता है।
\end{solution}

\begin{solution}{exo:b2:climate-variability:8}
(क) $(-1.93 - 1.11 - 1.98)/3 = -1.673$, लगभग $-1.67$। (ख) यह अभिलेख की सबसे निम्न शीत ऋतु है।
(ग) प्रबल ऋणात्मक NAO: दुर्बल आइसलैंडी निम्न और एज़ोर्स उच्च, दुर्बल जेट और सामान्य से
अधिक दक्षिण में तूफ़ान-पथ; उत्तरी यूरोप में ठंडी, शुष्क शीत ऋतु, दक्षिणी यूरोप में गीली,
और पूर्वी संयुक्त राज्य अमेरिका में ठंडी वायु के अधिक प्रकोप।
\end{solution}

\begin{solution}{exo:b2:climate-variability:9}
(क) $9.3 - 0.4 = \qty{8.9}{mm}$ प्रतिदिन, 2.5 से अधिक: क्षेत्र में। (ख) दक्षिणी गोलार्ध में
ग्रीष्म दिसंबर--फ़रवरी है: $6.1 - 4.2 = \qty{1.9}{mm}$ प्रतिदिन, 2.5 से कम: क्षेत्र से बाहर,
यद्यपि वहाँ क्षेत्र के अधिकांश स्थानों से अधिक वर्षा होती है। (ग) $1894/122 = \qty{15.5}{mm}$
प्रतिदिन; केरल केवल तभी बाहर होगा जब उसकी शीत ऋतु की दर $15.5 - 2.5 = \qty{13.0}{mm}$ प्रतिदिन से अधिक हो।
(घ) \omterm{def:b2:climate-variability:monsoon}{मानसून} एक ऋतुवार उलटाव है: पूरे वर्ष गीला रहने वाले गीले स्थान, जैसे दूसरा
केंद्र, में \omterm{def:b2:climate-variability:monsoon}{मानसून} नहीं होता, और छोटी, तीव्र वर्षा ऋतु वाले शुष्क स्थान में होता है।
\end{solution}

\begin{solution}{exo:b2:climate-variability:10}
(क) अवकलन करने पर $\ddot T = a\dot H = -abT$: कोणीय आवृत्ति $\omega = \sqrt{ab}$ और
आवर्तकाल $2\pi/\sqrt{ab}$ वाले आवर्ती दोलित्र का समीकरण। (ख) $2\pi/\sqrt{ab} = 4$ से $ab =
(2\pi/4)^2 = \qty{2.47}{yr^{-2}}$। (ग) $T = u\,\mathrm{e}^{-rt}$, $H = v\,\mathrm{e}^{-rt}$ के साथ,
$\dot T = (\dot u - ru)\,\mathrm{e}^{-rt} = (av - ru)\,\mathrm{e}^{-rt}$, अतः $\dot u = av$, और
इसी प्रकार $\dot v = -bu$: $(u, v)$ बिना अवमंदन के दोलन करता है, और $(T, H)$ वही दोलन है
जिसे $\mathrm{e}^{-rt}$ से गुणा किया गया है। एक आवर्तकाल में आयाम $\mathrm{e}^{-0.25 \times
4} = \mathrm{e}^{-1} = 0.37$ गुना रह जाता है। (घ) एक प्रणोदन: व्यापारिक पवनों के यादृच्छिक उतार-चढ़ाव, जो
तंत्र को बार-बार धक्का देते हैं, और \omterm{def:b2:climate-variability:bjerknes}{ब्येर्कनेस प्रतिपुष्टि}, जो अवमंदन को दुर्बल या
ऋणात्मक भी बना देती है।
\end{solution}

\begin{solution}{exo:b2:climate-variability:11}
(क) $A = R^2 \,\Delta\lambda\,(\sin\ang{5} - \sin(\ang{-5})) = (6.371 \times 10^6)^2 \times 0.8727
\times 0.1743 = \qty{6.17e12}{m^2}$। (ख) $Q = 1025 \times 3992 \times 100 \times 2 \times 6.17
\times 10^{12} = \qty{5.05e21}{J}$। (ग) छह महीने $182.6 \times 86\,400 = \qty{1.578e7}{s}$ हैं;
प्रति वर्ग मीटर, $8.18 \times 10^8/1.578 \times 10^7 = \qty{51.9}{W/m^2}$, आधे वर्ष तक
बना हुआ। (घ) स्वयं महासागर से: दुर्बल व्यापारिक पवनें पूर्व में कम ठंडा जल ऊपर लाती हैं, और पश्चिम का गर्म
जल, जो अधिक गहरी ताप-प्रवणता परत के ऊपर संचित है, पूर्व की ओर फैलता है; गर्म होने का बड़ा भाग
महासागर में पहले से मौजूद ऊष्मा से आ सकता है, जो प्राप्त होने के बजाय स्थानांतरित होती है।
\end{solution}

\begin{solution}{exo:b2:climate-variability:12}
(क) चक्र $0.3g$ लौटाता है, फिर $0.3g^2$, इत्यादि: कुल $0.3(1 + g + g^2 + \dots) =
0.3/(1 - g)$ है, एक गुणोत्तर श्रेणी जो $g < 1$ के लिए अभिसारी है: क्रमशः \qty{0.6}{K},
\qty{1.5}{K} और \qty{3.0}{K}, $g = 0.5$, 0.8 और 0.9 के लिए। (ख) श्रेणी अपसारी है: विसंगति असीम रूप से
बढ़ेगी; वृद्धि को उन प्रक्रियाओं द्वारा रोका जाना चाहिए जिन्हें चक्र छोड़ देता है, सतही ऊष्मा-प्रवाह द्वारा
अवमंदन और गर्म-जल भंडार का खाली होना, जिसका वर्णन पुनर्भरण दोलित्र करता है। (ग)
समुद्र-सतह जितनी गर्म होती है, उतनी ही आसानी से छोटा-सा गर्म होना गहरा संवहन और वर्षा शुरू कर देता है,
जिनका तापन पवनों को चलाता है; यही वह ऋतु है जिसमें ENSO के प्रभाव सबसे प्रबल होते हैं।
\end{solution}

\begin{solution}{pb:b2:climate-variability:1}
\textbf{1.} $(0.26 + 0.75 + 1.13)/3 = \qty{0.713}{\celsius}$।
\textbf{2.} MJJ $(0.75 + 1.13 + 1.52)/3 = 1.133$; JJA $(1.13 + 1.52 + 1.80)/3 = 1.483$; JAS $(1.52 +
1.80 + 2.06)/3 = 1.793$।
\textbf{3.} ASO $(1.80 + 2.06 + 2.26)/3 = 2.040$; SON $(2.06 + 2.26 + 2.42)/3 = 2.247$; OND $(2.26
+ 2.42 + 2.36)/3 = 2.347$; NDJ $(2.42 + 2.36 + 2.33)/3 = 2.370$।
\textbf{4.} एक दशमलव तक पूर्णांकित, 0.7, 1.1, 1.5, 1.8, 2.0, 2.2, 2.3 और 2.4: यानी तालिका।
\textbf{5.} एक अकेला महीना मौसम का रव लिए होता है; ENSO अनेक महीनों तक चलता है, और तीन का माध्य
धीमे संकेत को बनाए रखते हुए रव का बड़ा भाग हटा देता है।
\textbf{6.} JJA 1997 में, 1.483 पर, जो 1.5 तक पूर्णांकित होता है; वह NDJ में 2.4 पर शिखर पर पहुँची।
\textbf{7.} ASO 2018 से MJJ 2019 तक, और SON 2019 से FMA 2020 तक।
\textbf{8.} दस ऋतुएँ, और छह।
\textbf{9.} ASO 2020 में; 2020 के भीतर केवल चार ऋतुएँ, ASO से NDJ, $-0.5$ पर या उससे नीचे हैं, नियम से
एक कम, और DJF 2021 उसे पूरा करता है।
\textbf{10.} एक एल~नीनो, SON 2018 से MAM 2019, सात ऋतुएँ (AMJ 2019 गिरकर 0.37 पर आ जाता है);
ला~नीना AMJ 2020 ($-0.55$) में शुरू होती है, ONI द्वारा उसे रखे गए स्थान, ASO, से चार ऋतुएँ पहले।
\textbf{11.} नहीं: 2019 के अंत में सापेक्ष सूचकांक कभी 0.29 से ऊपर नहीं जाता। पूरा उष्णकटिबंधीय महासागर
तब अपनी आधार अवधि से लगभग आधा अंश अधिक गर्म था; नीनो-3.4 क्षेत्र
निरपेक्ष रूप से गर्म था, पर अपने आसपास के उष्णकटिबंधीय क्षेत्र से मुश्किल से अधिक गर्म।
\textbf{12.} ग्यारह।
\textbf{13.} 8, 7, 10, 4, 5, 6, 12, 5, 9 और 3 वर्ष।
\textbf{14.} सबसे छोटा 3 वर्ष (2023 से 2026), सबसे लंबा 12 (1997 से 2009)।
\textbf{15.} दोनों तरीकों से $69/10 = 6.9$ वर्ष: क्रमागत अंतरालों का योग पहली शुरुआत से
अंतिम तक का समय है, $2026 - 1957 = 69$ वर्ष।
\textbf{16.} वह पहले ही पाँच ऋतुओं, MAM से JAS, तक \qty{0.5}{\celsius} पर या उससे ऊपर रह चुकी है, और
2.2 पर, 1.5 से ऊपर, शिखर पर पहुँची है; उसकी शुरुआत, MAM 2026, क्रम की पहली ऋतु है, और बाद की ऋतुएँ
केवल क्रम को बढ़ा सकती हैं, उसकी शुरुआत को खिसका नहीं सकतीं।
\textbf{17.} नहीं, और होना भी आवश्यक नहीं: सभी तीव्रताओं की 25 घटनाएँ, जैसा FAQ कहता है, लगभग हर तीन वर्ष में आती हैं;
प्रबल घटनाएँ उनका एक उपसमुच्चय हैं, इसलिए दुर्लभतर हैं, और उनके अंतराल 3 से 12 वर्ष तक हैं।
\textbf{18.} $\tau = 1025 \times 3992 \times 50/16.1 = \qty{1.27e7}{s} = \qty{147}{\day}$।
\textbf{19.} मिश्रित परत: $1025 \times 3992 \times 50 = \qty{2.05e8}{J.m^{-2}.K^{-1}}$; वायु-स्तंभ:
$(10^5/9.81) \times 1004.5 = 10\,194 \times 1004.5 = \qty{1.02e7}{J.m^{-2}.K^{-1}}$, 20 गुना
कम।
\textbf{20.} $1.024 \times 10^7/(1025 \times 3992) = \qty{2.5}{m}$ समुद्री जल।
\textbf{21.} $1.024 \times 10^7/16.1 = \qty{6.36e5}{s} = 7.4$~दिन।
\textbf{22.} केवल वायु की विसंगति लगभग एक सप्ताह में अवमंदित हो जाती है, मिश्रित परत की लगभग
पाँच महीनों में, और ताप-प्रवणता परत के ऊपर की गर्म परत और भी गहरी है: केवल महासागर ऊष्मा को
इतनी देर तक रखता है कि विसंगति को एक ऋतु से अगली तक ले जा सके, जबकि वायुमंडल
दोनों को युग्मित करने वाली पवनें देता है।
\textbf{23.} 1950 के बाद से प्रबल एल~नीनो घटनाएँ औसतन हर \textbf{6.9 वर्ष} में शुरू हुई हैं,
लगभग सात।
\end{solution}
